\documentclass{article}
\usepackage{graphicx} % Required for inserting images
\usepackage{amsmath}
\usepackage[margin=.5in]{geometry}
\setlength{\parindent}{0pt}

\title{Homework 5}
\author{Quan Nguyen}
\date{March 2026}

\begin{document}
\maketitle

Homework code located at: https://github.com/nguyencquan/Markov2026

\section{Eigenvalues of Stochastic Matrices}

\textbf{a)} Show that p has at least one eigenvalue equal to 1.

If p is a stochasitc matrix size NxN, then all of the rows sum up to 1. We are trying to find some matrix
v such that $pv = v$. Expanding this matrix multiplication we get a specific index of v being
$$
v_{j} = \sum_{i=1}^{N}p_{j,i}v_{i}
$$
However we know that $\sum_{i=1}^n p_{j,i}=1,\forall j$, hence if we define every vector in v as 1, then we have
$$
v_{j} = \sum_{i=1}^{N}p_{j,i} = 1
$$
Therefore, for any stochastic matrix, an eigenvalue of 1 will have a corresponding matrix with
a column length N being composed of 1, as all rows in the stochastic matrix sum to 1.

\textbf{b) Show that all the other eigenvalues have magnitudes less then or equal to 1.}

The Gershgorin circle theorem states that, if $R_i = \sum_{i\neq j}|a_{ij}|$, where i is some row
of the matrix A, then all eigenvalues of A exists in at least one gershgorin disk denoted by
$(D(a_{ii},R_i))$

Looking at a stochastic matrix which has all positive entries and all rows sum up to one, we know for all i
$$1 = \sum_{j=1}^Na_{ij} = a_{ii} + \sum_{i\neq j}a_{ij} = a_{ii}+\sum_{i\neq j}|a_{ij}|$$

Hence the maximum upper bound for this Gershgorin disk is 1, and is always 1.

Furthermore, since $0\leq a_{ii} \leq 1$ we can calculate the lower bound from

\begin{align*}
    &1-a_{ii} = R_i\\
    &\implies 0\leq R_i \leq 1\\
    &\implies 1 \geq a_{ii} - R_i \geq -1
\end{align*}

From the first part we see the largest possible upper bound is 1 and lowest lower bound is $-1$
hence the absolute value of all eigenvalues must be equal to or less then 1.

\textbf{c) Show that if p1 and p2 are stochastic matrices of the same size, p1p2 is also stochastic}

In order to be a stochastic matrix, the sum of rows must equal 1 hence if we look at a single row.
\begin{align*}
    1 = \sum_{j=1}^\infty p1p2_{i,j} &= \sum_{k=1}^N\sum_{j=1}^N p1_{i,j}p2_{j,k}\\
    &= \sum_{j=1}^N\sum_{k=1}^N p1_{i,j}p2_{j,k}\\
    &= \sum_{j=1}^Np1_{i,j}\sum_{k=1}^N p1_{j,k}\\
    &= \sum_{j=1}^Np1_{i,j}=1\\
\end{align*}
However this applies to all rows so the product of these two matrices are stochastic.
\end{document}